% Generated draft: requires mathematical and source review.
\documentclass[11pt]{article}
\usepackage[margin=25mm]{geometry}
\usepackage{fontspec}
\setmainfont{Latin Modern Roman}
\usepackage{amsmath,amssymb,mathtools,unicode-math}
\setmathfont{Latin Modern Math}
\usepackage{enumitem,needspace,xurl,hyperref}
\setlength{\parindent}{0pt}
\setlength{\parskip}{5pt}
\setlength{\emergencystretch}{3em}
\newcommand{\N}{\mathbb{N}}
\newcommand{\Z}{\mathbb{Z}}
\newcommand{\Q}{\mathbb{Q}}
\newcommand{\R}{\mathbb{R}}
\newcommand{\C}{\mathbb{C}}
\newcommand{\F}{\mathbb{F}}
\newcommand{\A}{\mathbb{A}}
\newcommand{\PP}{\mathbb{P}}
\newcommand{\E}{\mathbb{E}}
\newcommand{\eps}{\varepsilon}
\newcommand{\dd}{\,\mathrm{d}}
\begin{document}
\section*{UnsolvedMath ID: 30003260}
\textbf{Source problem code:} catalogue entry
\subsection{Definitions and mathematical statement}
Let \(p\) be a prime, let \(t\) be transcendental over \(\F_p\), and put \(K=\F_p(t)\). For \(g\in K\), let \(\F_p(g)\) denote the smallest subfield of \(K\) containing \(\F_p\) and \(g\). This definition includes constant \(g\), for which \(\F_p(g)=\F_p\). Define
\[
S=\{(f,g)\in K^2:f\in\F_p(g)\}.
\]
A subset \(D\subseteq K^n\) is \emph{Diophantine in the parameter-free ring language} if there exist integers \(m\geq0\), \(r\geq1\), and polynomials
\[
P_1,\ldots,P_r\in\Z[X_1,\ldots,X_n,Y_1,\ldots,Y_m]
\]
such that, for every \(\mathbf x\in K^n\),
\[
\mathbf x\in D\quad\Longleftrightarrow\quad
\exists\mathbf y\in K^m\ \bigwedge_{i=1}^{r}P_i(\mathbf x,\mathbf y)=0.
\]
Integer coefficients are interpreted through the canonical map \(\Z\to K\); in particular, no parameter naming \(t\) is allowed. Diophantine definability \emph{with parameters in \(K\)} is defined in the same way, allowing the finitely many defining polynomials to have coefficients in \(K\).

\textbf{Question.} Is \(S\subseteq K^2\) Diophantine in the parameter-free ring language? Explicitly, do there exist \(m\geq0\), \(r\geq1\), and \(P_1,\ldots,P_r\in\Z[X,Z,Y_1,\ldots,Y_m]\) such that
\[
\forall f,g\in K,\qquad
f\in\F_p(g)\quad\Longleftrightarrow\quad
\exists\mathbf y\in K^m\ \bigwedge_{i=1}^{r}P_i(f,g,\mathbf y)=0?
\]
The source also asks, in particular, whether the subfield \(\F_p(t^2)\subseteq K\) is Diophantine with parameters in \(K\). That is, do there exist \(a\geq0\), \(b\geq1\), and \(Q_1,\ldots,Q_b\in K[X,Y_1,\ldots,Y_a]\) such that
\[
\forall f\in K,\qquad
f\in\F_p(t^2)\quad\Longleftrightarrow\quad
\exists\mathbf y\in K^a\ \bigwedge_{j=1}^{b}Q_j(f,\mathbf y)=0?
\]
A positive answer for \(S\) implies a positive answer to this latter question by specializing \(g=t^2\), which is then an allowed parameter. The supplied record gives no restriction on the prime \(p\), and does not specify whether the question concerns a particular fixed prime, every prime separately, or a single formula uniform in \(p\). The displayed questions describe definability in each fixed field \(\F_p(t)\); the intended quantification over primes remains unresolved.
\subsection{Short English statement}
In the rational function field \(\F_p(t)\), can membership of \(f\) in the subfield generated by \(g\) be expressed by existential polynomial equations without parameters? The source also asks whether membership in \(\F_p(t^2)\) can be expressed by existential polynomial equations allowing parameters from \(\F_p(t)\). The intended scope over primes is unspecified.
\subsection{Sources}
\begin{itemize}
\item Supplied local UnsolvedMath registry record, permanent ID \(30003260\), title \emph{Diophantine Definability of Rational Subfields}, in the Oberwolfach Reports Open Problems collection; exact source problem code \texttt{catalogue entry}.
\item Supplied report citation: \emph{Definability and Decidability Problems in Number Theory} (2017), Oberwolfach Reports, pp.~2793--2866, DOI \texttt{10.4171/owr/2016/49}. Supplied literature-source label: ``OWR 49 report.'' \url{https://doi.org/10.4171/owr/2016/49}
\item Supplied literature-source label: ``EMS report text.'' \url{https://ems.press/content/serial-article-files/46652}
\item The registry records proposed year 2016 and supplies no individual proposer attribution. Its report citation uses 2017; these supplied dates are retained without resolving their relationship.
\item Literature assessment dated exactly \texttt{2026-08-21}: the supplied curation labels the problem open and reports that no unconditional resolution was found in the subsequent literature checked. This is supplied dated curation, not an independently verified assessment. The background explicitly requests independent verification of time-sensitive status and cited claims.
\item The supplied statement status is \texttt{corrected\_verified}; its verification note reports a manual semantic, typographic, and LaTeX review preserving the mathematical target. This is registry metadata and does not constitute independent mathematical verification of this draft.
\end{itemize}
\end{document}
